% !TEX root = ../../main.tex
% C5.3 — Thiết kế cơ sở dữ liệu (RÚT GỌN 2026-09-29: bỏ 6 bảng từ điển dữ liệu và bảng các bảng vận hành; bản cũ ở report/archive/)
% 5.3.1 khớp backend/src/person_search/storage/postgres/models/*.py và 13 migration (backend/migrations/versions). 12 bảng.
% Bản số trong ERD lấy từ nullable của FK (xem tools/h8_erd_drawio.py). worker_heartbeats không có FK -> không vẽ trong ERD.
% camera "loại khỏi vận hành" = status RETIRED + ai_enabled=false (services/cameras.py); INACTIVE có trong enum nhưng chưa dùng.
% Bất biến areas.code, cameras.area_id và chuyển trạng thái person_tracks được kiểm tra ở tầng ứng dụng (sự kiện before_update của ORM).
% 5.3.2 khớp storage/contracts.py, storage/milvus/vectors.py (7 trường, HNSW/IP M=16 efConstruction=128 ef=64, Strong, upsert).
% 5.3.3 khớp storage/minio/frames.py (JPEG/PNG <= 20 MB, SHA-256, put idempotent), workers/publication.py (JPEG 90), infra/minio (bucket riêng tư).
% 5.3.4 khớp services/track_ingestion.py (PENDING + outbox cùng giao dịch -> MinIO -> Milvus + xác minh -> READY; thử lại 2 s nhân đôi <= 300 s,
% tối đa 5 lần) và storage/maintenance_cli.py (retry-outbox, reconcile mặc định chỉ báo cáo, requeue-track, reindex; chạy thủ công). Hình H6.
\section{Thiết kế cơ sở dữ liệu}
\label{sec:5.3}

\subsection{Cơ sở dữ liệu quan hệ}
\label{sec:5.3.1}

PostgreSQL lưu dữ liệu nghiệp vụ và là nơi quyết định về quyền và trạng thái. 12 bảng của nó chia thành năm nhóm (Bảng~\ref{tab:db-tables}).

\begin{table}[H]
\centering
\caption{Các bảng trong cơ sở dữ liệu quan hệ}
\label{tab:db-tables}
\begin{tabularx}{\textwidth}{|>{\raggedright\arraybackslash}p{3.0cm}|>{\raggedright\arraybackslash}p{5.4cm}|L|}
\hline
\textbf{Nhóm} & \textbf{Bảng} & \textbf{Vai trò} \\
\hline
Tổ chức và phân quyền & \texttt{areas}, \texttt{users}, \path{auth_sessions} & Khu vực giám sát, tài khoản người dùng, phiên đăng nhập \\
\hline
Camera và xử lý AI & \texttt{cameras}, \path{ai_config_versions}, \path{processing_jobs}, \path{worker_heartbeats} & Camera logic, các phiên bản cấu hình Detector-Tracker-bộ mã hóa, hàng đợi công việc, trạng thái tiến trình xử lý nền \\
\hline
Dữ liệu tìm kiếm & \path{person_tracks}, \path{storage_outbox_events} & Thông tin mô tả các lần xuất hiện; sự kiện đồng bộ với hai kho còn lại (Mục~\ref{sec:5.3.4}) \\
\hline
Nghiệp vụ vụ việc & \texttt{cases}, \path{case_results} & Vụ việc và các kết quả đã lưu \\
\hline
Truy vết & \path{audit_logs} & Nhật ký hệ thống, chỉ được thêm, không sửa \\
\hline
\end{tabularx}
\end{table}

Hình~\ref{fig:erd} là sơ đồ thực thể - liên kết của 11 bảng có khóa ngoại theo ký hiệu chân quạ, gồm khóa và các thuộc tính chính; bảng \path{worker_heartbeats} không có khóa ngoại nên không được vẽ.

\afterpage{%
\begin{landscape}
\begin{figure}[p]
\centering
\includegraphics[width=\linewidth,height=0.86\textheight,keepaspectratio]{H8-erd.png}
\caption{Sơ đồ thực thể - liên kết của cơ sở dữ liệu quan hệ}
\label{fig:erd}
\end{figure}
\end{landscape}
}

\paragraph{Các quan hệ.} Phần lớn khóa ngoại là bắt buộc: camera có một khu vực, vụ việc có một người phụ trách, lần xuất hiện có một camera, một công việc và một phiên bản cấu hình AI. Ba khóa ngoại được rỗng: chỉ Giám sát viên có khu vực, chỉ công việc từ tệp tải lên có người yêu cầu, và bản ghi nhật ký như một lỗi kỹ thuật có thể không gắn với người dùng. Giữa \texttt{cases} và \path{person_tracks} là quan hệ nhiều - nhiều qua bảng \path{case_results}. Mỗi dòng của bảng này có khóa chính riêng, lưu bản chụp thông tin tại thời điểm lưu, và \emph{không} có ràng buộc duy nhất trên cặp vụ việc - lần xuất hiện, vì mỗi lần lưu tạo một mục mới.

\paragraph{Các ràng buộc toàn vẹn.} Nhiều quy tắc nghiệp vụ do chính cơ sở dữ liệu ràng buộc, nên lỗi ở tầng ứng dụng không ghi được dữ liệu sai:
\begin{itemize}
    \item Tài khoản Giám sát viên luôn có khu vực, hai vai trò còn lại thì không; vụ việc có thời điểm hoàn thành khi và chỉ khi ở trạng thái \texttt{CLOSED}; chỉ một phiên bản cấu hình AI ở trạng thái \texttt{ACTIVE}; mã khu vực và khu vực của camera không đổi sau khi tạo.
    \item Hầu hết khóa ngoại dùng quy tắc RESTRICT, nên dữ liệu còn được tham chiếu không thể bị xóa; tài khoản và camera không bị xóa dòng mà được chuyển trạng thái (\texttt{INACTIVE}, \texttt{RETIRED}). Chỉ phiên đăng nhập và mục kết quả được xóa theo bảng cha.
    \item Một lần xuất hiện chỉ có thể ở trạng thái \texttt{READY} khi đã có đủ khóa đối tượng, mã băm và kích thước khung hình, và thời điểm ghi vector. Không bảng nào lưu điểm phù hợp.
    \item Các bảng \texttt{users}, \texttt{cameras} và \texttt{cases} có cột số phiên bản để chống ghi đè khi cập nhật đồng thời (Mục~\ref{sec:5.2.4}); cặp người yêu cầu - khóa chống trùng của \path{processing_jobs} là duy nhất, nên một yêu cầu tải lên gửi lặp không tạo thêm công việc.
\end{itemize}
Các chỉ mục phục vụ truy vấn thường gặp như lần xuất hiện theo camera và thời điểm, vụ việc theo người phụ trách, nhật ký theo người thực hiện.

\subsection{Lưu trữ vector đặc trưng}
\label{sec:5.3.2}

Milvus chỉ tìm các vector gần nhất với truy vấn trong một phạm vi đã lọc; mọi thứ khác, kể cả quyền, nằm ở PostgreSQL.

\paragraph{Mỗi phiên bản bộ mã hóa một collection.} Mỗi phiên bản bộ mã hóa có một \emph{collection} (tương tự bảng) riêng trong Milvus, ví dụ \path{person_track_embeddings_rasa_cuhk_pedes_v1}, vì vector của hai phiên bản không so sánh được với nhau. Một bí danh cố định trỏ tới collection đang dùng, và số chiều được kiểm tra khi tiến trình nền khởi động. Mỗi lần xuất hiện ứng với một bản ghi gồm các trường ở Bảng~\ref{tab:milvus-schema}.

\begin{table}[H]
\centering
\caption{Lược đồ collection lưu vector đặc trưng trong Milvus}
\label{tab:milvus-schema}
\begin{tabularx}{\textwidth}{|>{\raggedright\arraybackslash}p{4.2cm}|>{\raggedright\arraybackslash}p{3.7cm}|L|}
\hline
\textbf{Trường} & \textbf{Kiểu} & \textbf{Ý nghĩa} \\
\hline
\path{track_id} & VARCHAR(36), khóa chính & Mã lần xuất hiện, trùng với khóa chính của \path{person_tracks} \\
\hline
\path{embedding} & FLOAT\_VECTOR, 256 chiều & Vector đặc trưng, đã chuẩn hóa L2 \\
\hline
\path{area_id}, \path{camera_id} & VARCHAR(36) & Khu vực và camera, dùng để lọc theo phạm vi \\
\hline
\path{appeared_at_epoch} & INT64 & Thời điểm xuất hiện (giây Unix), dùng để lọc theo thời gian \\
\hline
\path{encoder_version} & VARCHAR(64) & Phiên bản bộ mã hóa đã tạo vector \\
\hline
\path{index_status} & VARCHAR(16) & Trạng thái sẵn sàng; chỉ bản ghi \texttt{READY} được tìm \\
\hline
\end{tabularx}
\end{table}

Khu vực, camera và thời điểm được sao từ PostgreSQL để lọc trước (Mục~\ref{sec:3.6.3}); khu vực của camera không đổi nên bản sao không thể lệch. Trạng thái vận hành thì thay đổi nên không sao chép: mỗi lượt tìm đưa danh sách camera đang vận hành vào điều kiện lọc, nên camera bị loại được ẩn ngay mà không phải sửa Milvus.

\paragraph{Chỉ mục và thao tác.} Trường vector được đánh chỉ mục HNSW với độ đo tích vô hướng, $M = 16$, \textit{ef\-Con\-struc\-tion}~$= 128$, và $\mathit{ef} = 64$ khi tìm (Mục~\ref{sec:3.6.4}). Vector được ghi bằng \emph{upsert} theo mã lần xuất hiện nên thử lại không tạo bản ghi trùng, sau khi kiểm tra số chiều, giá trị và độ dài bằng 1; mức nhất quán mạnh (strong consistency) giúp vector vừa ghi có ngay trong lần tìm tiếp theo.

\subsection{Lưu trữ khung hình}
\label{sec:5.3.3}

Mỗi lần xuất hiện có một tệp trong MinIO là khung hình toàn cảnh đại diện, JPEG chất lượng 90; PostgreSQL lưu khóa đối tượng, mã băm và kích thước. Ảnh người không được lưu mà dựng từ khung hình và khung bao khi cần (Mục~\ref{sec:6.3.3}), nên một tệp phục vụ cả hai cách xem, cách cắt đổi được mà không tạo lại dữ liệu, và không có tệp ảnh người rời. Video gốc không được lưu.

Khóa của mỗi tệp được sinh xác định từ thông tin của lần xuất hiện, theo mẫu
\begin{center}
\texttt{tracks/v1/}\allowbreak \textit{mã camera}\texttt{/}\allowbreak \textit{năm}\texttt{/}\allowbreak \textit{tháng}\texttt{/}\allowbreak \textit{ngày}\texttt{/}\allowbreak \textit{mã lần xuất hiện}\texttt{/}\allowbreak \texttt{representative.jpg}
\end{center}
nên thử lại chỉ ghi đè đúng tệp đó; tiền tố có số phiên bản cho phép đổi cách tổ chức về sau, và nhóm theo camera và ngày giúp dọn dẹp dễ.

Tệp được kiểm tra định dạng, dung lượng (không quá 20~MB) và kích thước, và mã băm SHA-256 lưu ở cả hai kho: trùng khóa cùng mã băm thì dùng lại tệp, khác mã băm là xung đột, và khi đọc mã băm được tính lại để phát hiện tệp hỏng. Bucket là riêng tư, truy cập bằng tài khoản chỉ có quyền trên nó, và không lộ ra trình duyệt.

\subsection{Đảm bảo nhất quán giữa các kho dữ liệu}
\label{sec:5.3.4}

Một lần xuất hiện được lưu thành ba phần ở ba kho không có giao dịch chung, nên việc ghi có thể dừng giữa chừng hoặc một kho có thể tạm không truy cập được. Hệ thống bảo đảm dữ liệu dở dang không xuất hiện trong kết quả, và có thể được hoàn tất hoặc phát hiện về sau.

\paragraph{Trạng thái của lần xuất hiện.} Mỗi lần xuất hiện có một trạng thái đồng bộ với các chuyển trạng thái ở Hình~\ref{fig:track-state}; chỉ lần xuất hiện \emph{Sẵn sàng} mới được tìm, và các chuyển trạng thái khác bị từ chối.

\begin{figure}[!htbp]
\centering
\includegraphics[width=\textwidth]{H6-trang-thai-lan-xuat-hien.png}
\caption{Sơ đồ trạng thái của một lần xuất hiện}
\label{fig:track-state}
\end{figure}

\paragraph{Thứ tự ghi.} Việc ghi một lần xuất hiện theo thứ tự cố định:
\begin{enumerate}
    \item Trong \emph{một giao dịch} của PostgreSQL, tạo bản ghi ở trạng thái \emph{Chờ} và một sự kiện trong bảng \path{storage_outbox_events} lưu đủ những gì cần ghi sang hai kho còn lại, kể cả vector. Đây là mẫu thiết kế \emph{transactional outbox}: ý định ghi sang hệ thống khác được lưu cùng giao dịch với dữ liệu chính.
    \item Ghi khung hình vào MinIO theo khóa xác định.
    \item Ghi vector vào Milvus, rồi đọc lại để xác nhận bản ghi đã tồn tại với đúng khu vực và camera.
    \item Trong một giao dịch mới, chuyển bản ghi sang \emph{Sẵn sàng}, lưu mã băm, kích thước khung hình và thời điểm ghi vector, và đánh dấu sự kiện outbox đã hoàn tất.
\end{enumerate}
Ghi PostgreSQL \emph{trước} nên mọi khung hình và vector đều đối chiếu được với một bản ghi, phần còn lại là dữ liệu mồ côi; đặt \emph{Sẵn sàng} \emph{cuối cùng} nên lần xuất hiện chỉ tìm được khi đủ dữ liệu. Bước 2 và 3 lặp lại an toàn nhờ khóa xác định và upsert.

\paragraph{Xử lý lỗi và công cụ bảo trì.} Lỗi tạm thời ở bước 2 hoặc 3 giữ bản ghi ở \emph{Chờ} và hẹn thử lại, từ 2 giây tăng gấp đôi tới tối đa 300 giây. Sau năm lần, hoặc khi gặp lỗi không phục hồi như xung đột mã băm, bản ghi chuyển sang \emph{Thất bại} và được ghi nhật ký. Vì sự kiện outbox chứa đủ dữ liệu, việc thử lại không phải chạy lại phân tích. Người vận hành có thêm một công cụ dòng lệnh để hoàn tất các lần xuất hiện còn dở, đối soát ba kho nhằm tìm bản ghi chờ quá lâu, bản ghi sẵn sàng nhưng thiếu khung hình hoặc vector, và dữ liệu mồ côi (mặc định chỉ báo cáo; khi được yêu cầu thì cách ly bản ghi hỏng và xóa dữ liệu mồ côi), đưa một lần xuất hiện thất bại về trạng thái chờ, và dựng lại dữ liệu Milvus từ các sự kiện outbox. Các lệnh này hiện được chạy thủ công.

\paragraph{Lớp bảo vệ khi tìm kiếm.} Kể cả khi các kho chưa khớp, người dùng không thấy dữ liệu sai: mỗi kết quả Milvus đều được đối chiếu với PostgreSQL, và kết quả không tồn tại hoặc ngoài phạm vi bị loại (Mục~\ref{sec:3.6.4}).
